\section{Discussion}
\label{sec:discussion}

The empirical findings across Stages S0, S1, and S2 demonstrate that the aggregate output--capital relation does not constitute an insulated engineering benchmark. Cointegration fails entirely in the standalone bivariate system over the post-war period, and within the tested full-sample system grid, long-run stability becomes recoverable only when the rate of exploitation enters the state vector alongside historical step controls. This empirical contrast motivates the central conceptual interpretation of this paper: productive capacity is an institutionally conditioned outcome of capitalist accumulation.

Crucially, ``institutional settlement'' must be understood as an analytical interpretation of the econometric results rather than an estimated parameter. The VECM does not directly measure class struggle or institutional power. Instead, the institutional settlement framework explains why the bivariate relation fails in isolation and why distribution enters the cointegrating space with such dominant statistical precision ($\hat{\beta}_e = 20.02, t = 7.90$, compared to the imprecisely identified capital elasticity $\hat{\theta} = 0.73, \text{SE} = 4.85$). In capitalist economies, the conversion of a given stock of physical equipment into potential output is not governed solely by technical engineering coefficients. It depends on the organization of the labor process, machine speeds, shift patterns, and managerial authority on the shop floor. These conditions are mediated by the historical compromises and power balances between capital and labor that determine the functional distribution of income. When distribution is omitted, the tested bivariate systems exhibit persistent drift in the output--capital relation. The recovery of cointegration in specifications including $\ln e_t$ is consistent with distributionally mediated adjustments in technique and work organization, but the VECM does not identify that mechanism causally.

While Stage S1 provides the primary statistical evidence for sub-unitary elasticities across retained no-trend specifications in the bounds-passing ARDL grid ($\hat{\theta} \in [0.65, 0.95]$), the focal point estimate in Stage S2 ($\hat{\theta} \approx 0.73$) also sits within the structural overaccumulation interval ($\theta < 1.0$), albeit with wide normalized estimation uncertainty. If interpreted through this point estimate, the profit rate decomposition,
\begin{equation}
r_t \equiv \frac{\Pi_t}{K_t} = \left(\frac{\Pi_t}{Y_t}\right) \left(\frac{Y_t}{Y^p_t}\right) \left(\frac{Y^p_t}{K_t}\right) = \pi_t \, \mu_t \, \left(\frac{Y^p_t}{K_t}\right),
\end{equation}
carries significant implications for classical and Marxian crisis theory. The maximum potential profit rate is governed by the capital-capacity ratio $Y^p_t / K_t$. When $\theta = 1.0$, capital accumulation expands capacity proportionally, preserving the capacity-capital ratio in the absence of technical change. When $\theta < 1.0$, however, the corporate sector operates under conditions of structural overaccumulation: capital equipment accumulates faster than potential output expands, placing downward pressure on $Y^p_t / K_t$ and, conditional on distribution, on potential profitability. A higher profit share ($\pi_t$) can partly offset this pressure in the realized profit rate, whereas utilization below normal capacity ($\mu_t < 1$) further depresses realized profitability.

The dynamic adjustment structure of the focal VECM also speaks to the long-standing debate between Sraffian and Neo-Kaleckian macroeconomics. In the focal system, the estimated error-correction loading for capital is statistically indistinguishable from zero ($\hat{\alpha}_k = 0.0003, t=0.92$), while the primary burden of error correction falls on the rate of exploitation ($\hat{\alpha}_e = -0.019, t=-4.25$). The absence of statistically detectable error-correction adjustment in capital accumulation is suggestively aligned with classical-Sraffian models in which long-run investment is autonomous, rather than with Neo-Kaleckian models where investment responds endogenously to capacity utilization gaps. Nevertheless, this point-estimate asymmetry should not be overstated as a definitive statistical refutation of Neo-Kaleckian closure without a joint likelihood-ratio restriction test. Moreover, insofar as sub-unitary elasticity estimates ($\hat{\theta} < 1.0$) decouple capacity from capital, the long-run dynamics of accumulation cannot be evaluated under the balanced-growth knife-edge assumption that both traditions routinely share.

Several methodological boundaries qualify these findings. First, the empirical analysis relies on single-sector corporate data, abstracting from inter-sectoral transfers, department-level disproportionalities, and uneven development across branches of production. Second, the transformation elasticity $\theta$ is estimated as a fixed parameter over the 1947--2011 period, averaging across distinct institutional regimes (such as the post-war Golden Age and the neoliberal transition). Third, the bivariate failure reflects sample sensitivity: while bivariate systems yield admissible cointegrating relations over the 1947--2007 sub-sample, cointegration is lost over the full 1947--2011 sample, coinciding with the acute contraction of the Great Recession where value added dropped sharply relative to the corporate capital stock.

\section{Conclusion}
\label{sec:conclusion}

This paper critically evaluated Anwar Shaikh's (2016) capacity utilization measure, investigating whether the underlying aggregate output--capital relation provides a stable technical benchmark for macroeconomic analysis. Re-examining this relation through an econometric replication across three stages yields three core findings that establish the evidentiary ordering of the paper.

First, at the single-equation level (Stage S0), Shaikh's baseline ARDL(2,4) specification is approximately reproducible within canonical data bounds ($\hat{\theta} \approx 0.72$ versus Shaikh's published 0.66). Second, testing this relationship across an exhaustive lattice of 500 ARDL specifications (Stage S1) reveals disciplined non-uniqueness: the estimated transformation elasticity ranges across the informational frontier from 0.65 to 0.95 among retained no-trend specifications (reaching 2.20 in trend-containing models). The recovered capacity ceiling is sensitive to lag profiles and deterministic trend assumptions, demonstrating that single-equation models cannot identify a unique technical parameter.

Third, system-level estimation (Stage S2) demonstrates that across 48 attempted bivariate VECM specifications, 36 are successfully estimated, and none of those 36 identifies an admissible cointegrating relation over the 1947--2011 sample. Within the tested full-sample system grid, cointegration becomes recoverable only when the rate of exploitation enters the state vector alongside historical step controls. In the focal trivariate VECM, the rate of exploitation enters with dominant statistical precision ($\hat{\beta}_e = 20.02, t = 7.90$), while the transformation elasticity point estimate ($\hat{\theta} \approx 0.73$) sits within the structural overaccumulation interval ($\theta < 1.0$) alongside substantial estimation uncertainty.

The overarching conclusion is that aggregate productive capacity is not an insulated engineering constant. Rather, the output--capital relation is distributionally conditioned and socially structured. Macroeconomic models that treat the capacity ceiling as an autonomous physical limit risk obscuring the institutional arrangements and distributive conflicts that govern how productive equipment is converted into potential output in capitalist economies.
